\subsection{DEFINITION OF AN INTEGRAL OVER A NON-COMPACT INTERVAL}

Let I be a compact interval \([a, b]\) in the extended line \(\overline{R}\) (a and b may be infinite); let f be a function defined on {]}a, b{[}, taking its values in a complete normed space E over R. Generalizing def. 1 of II, p.~51, we shall say that a function g, defined on \([a, b]\) with values in E, is a primitive of f if it is continuous on \([a, b]\) (and in particular at the endpoints a and b) and admits a derivative equal to \(\mathbf{f}(x)\) at all the points of the complement with respect to {]}a, b{[} of a countable subset of this interval.

We shall restrict ourselves to the following case: there exists a finite strictly increasing sequence \((c_i)_{0\leqslant i\leqslant n}\) of points of \(\mathrm{I} = [a,b]\) , such that \(c_0 = a\) , \(c_n = b\) , and such that \(\mathbf{f}\) is regulated on each of the open intervals \(]c_i\) , \(c_{i + 1}[\) , though not necessarily regulated on every open interval containing at least one point \(c_i\) interior to \(\mathrm{I}\) ; such a function will be called piecewise regulated on \(]a\) , \(b[\) . We remark that a regulated function on \(]a\) , \(b[\) is piecewise regulated (taking \(n = 1\) in the preceding definition).

If \(\mathbf{f}\) admits a primitive \(\mathbf{g}\) (in the sense made precise above), and if \(c\) is a point of the interval \([c_i, c_{i+1}[(0 \leqslant i \leqslant n-1)\) , one has, by hypothesis, for each \(x\) in this interval, that \(\mathbf{g}(x) - \mathbf{g}(c) = \int_c^x \mathbf{f}(t) dt\) ; since \(\mathbf{g}\) is continuous on I by hypothesis, one sees that \(\int_c^x \mathbf{f}(t) dt\) tends to a limit in E when \(x\) tends to \(c_i\) from the right and when \(x\) tends to \(c_{i+1}\) from the left. Conversely, suppose that these conditions are satisfied for all \(t\) , and let \(\mathbf{g}_i\) be a primitive of \(\mathbf{f}\) on the interval \([c_i, c_{i+1}[(0 \leqslant i \leqslant n-1)\) ; we note immediately that the function \(\mathbf{g}\) , defined on the complement with respect to I of the set of \(c_i\) , by the condition that it be equal to \(\mathbf{g}_i(x) + \sum_{k=1}^i (\mathbf{g}_{k-1}(c_k-) - \mathbf{g}_k(c_k+))\) on \([c_i, c_{i+1}[\) for \(0 \leqslant i \leqslant n-1\) , is continuous at every point of I distinct from the \(c_i\) and admits a limit at each of these points; it can therefore be extended by continuity to each of these points, and the extended function is evidently a primitive of \(\mathbf{f}\) on I. It is clear, moreover, that every other primitive of \(\mathbf{f}\) is of the form \(\mathbf{g} + \mathbf{a}\) ( \(\mathbf{a}\) an element of E).

DEFINITION 1. One says that a vector function \(\mathbf{f}\) , piecewise regulated on an interval \(]a\) , \(b[\) of \(\overline{\mathbf{R}}\) , admits an integral on this interval if \(\mathbf{f}\) admits a primitive on \([a, b]\) ; if \(\mathbf{g}\) is any one of the primitives of \(\mathbf{f}\) on \([a, b]\) , and \(x_0\) and \(x\) are any two points of \([a, b]\) , one calls the element \(\mathbf{g}(x) - \mathbf{g}(x_0)\) the integral of \(\mathbf{f}\) from \(x_0\) to \(x\) , and one denotes it by \(\int_{x_0}^x \mathbf{f}(t) dt\) .

This concept clearly agrees with that defined when the interval \([x_0, x]\) contains none of the points \(c_i\) .

The remarks which precede def. 1 show that for \(\mathbf{f}\) to have an integral on \(]a\) , \(b[\) it is necessary and sufficient that its restriction to each of the intervals \([c_i, c_{i+1}]\) should admit an integral over this interval. In other words, one reduces to the case where \(\mathbf{f}\) is regulated on a non-compact interval \(\mathrm{I} \subset \mathbf{R}\) , with endpoints \(a, b\) ( \(a < b\) ), and where: \(1^\circ\) either one of the numbers \(a, b\) (at least) is infinite; \(2^\circ\) or \(\mathbf{f}\) is not regulated on a compact interval containing at least one of the points \(a, b\) (these two hypotheses not being mutually exclusive). For \(\mathbf{f}\) to have an integral over \(\mathrm{I}\) it is necessary and sufficient that the integral \(\int_x^y \mathbf{f}(t) dt\) approaches a limit when the point \((x, y)\) approaches \((a, b) \in \overline{\mathbf{R}}^2\) while remaining in \(\mathrm{I} \times \mathrm{I}\) : and this limit is no other than \(\int_a^b \mathbf{f}(t) dt\) according to def. 1. By an abuse of language, instead of saying that \(\mathbf{f}\) has an integral over \(\mathrm{I}\) , one says that the integral \(\int_a^b \mathbf{f}(t) dt\) is convergent (or converges).

Examples. 1) The integral \(\int_{1}^{+\infty} dt / t^{2}\) is convergent and equal to 1, for

\[
\int_ {1} ^ {x} \frac {d t}{t ^ {2}} = 1 - \frac {1}{x}.
\]

\begin{enumerate}
\def\labelenumi{\arabic{enumi})}
\setcounter{enumi}{1}
\tightlist
\item
  The integral \(\int_0^1 dt / \sqrt{t}\) is convergent and equal to 2, for
\end{enumerate}

\[
\int_ {x} ^ {1} \frac {d t}{\sqrt {t}} = 2 (1 - \sqrt {x}) \quad \text { for } \quad x > 0.
\]

\begin{enumerate}
\def\labelenumi{\arabic{enumi})}
\setcounter{enumi}{2}
\tightlist
\item
  Let \((\mathbf{u}_n)_{n\geqslant 1}\) be an infinite sequence of points of E, and let \(\mathbf{f}\) be the step function defined on the interval {[}1, \(+\infty[\) by the conditions: \(\mathbf{f}(x) = \mathbf{u}_n\) for \(n\leqslant x < n + 1\) . Then for the integral \(\int_1^{+\infty}\mathbf{f}(t)dt\) to be convergent it is necessary and sufficient that the series with general term \(\mathbf{u}_n\) be convergent in E; indeed, one has
\end{enumerate}

\[
\int_ {1} ^ {n} \mathbf {f} (t) d t = \sum_ {p = 1} ^ {n - 1} \mathbf {u} _ {p},
\]

so the condition is necessary; conversely, if the series with general term \(\mathbf{u}_n\) converges in E, then \(\lim_{n\to \infty}\mathbf{u}_n = 0\) ; now, if \(n\leqslant x\leqslant n + 1\) , one has \(\int_1^n\mathbf{f}(t)dt = \sum_{p = 1}^{n - 1}\mathbf{u}_p +\) \(\mathbf{u}_n(x - n)\) , so this integral has the limit \(\sum_{n = 1}^{\infty}\mathbf{u}_n\) when \(x\) tends to \(+\infty\) .

It is immediate that if a piecewise regulated function \(\mathbf{f}\) admits an integral over I then the formulae (4) to (9) of II, p.~59 remain valid. Similarly, formula (10) of II, p.~59 extends in the following manner: \(\mathbf{f}\) and \(\mathbf{g}\) are assumed to be primitives of the regulated functions \(\mathbf{f}'\) and \(\mathbf{g}'\) on \(]a\) , \(b[\) , and one denotes by \([\mathbf{f}.\mathbf{g}]|_a^b\) the limit (if it exists) of \([\mathbf{f}.\mathbf{g}]|_x^y\) as \((x, y)\) tends to \((a, b)\) (with \(a < x \leqslant y < b\) ); then, if two of the (three) expressions \([\mathbf{f}.\mathbf{g}]|_a^b\) , \(\int_a^b [\mathbf{f}(t).\mathbf{g}'(t)] dt\) , and \(\int_a^b [\mathbf{f}'(t).\mathbf{g}(t)] dt\) have a meaning, then so has the third, and the formula (10) of II, p.~59 is valid.

Finally, let \(f\) be a real function which is defined and continuous on \(I = ]a\) , \(b[\) , and is the primitive of a regulated function \(f'\) on \(]a\) , \(b[\) ; let on the other hand \(\mathbf{g}\) be a continuous vector function on an open interval \(J\) containing \(f(I)\) ; if the function \(\mathbf{g}(f(x))f'(x)\) admits an integral over \(I\) , and if \(f\) tends to a limit (finite or not) at the points \(a\) and \(b\) , then \(\mathbf{g}\) admits an integral from \(f(a+)\) to \(f(b-)\) , and one has the formula

\[
\int_ {a} ^ {b} \mathbf {g} (f (t)) f ^ {\prime} (t) d t = \int_ {f (a +)} ^ {f (b -)} \mathbf {g} (u) d u.\tag{1}
\]

Indeed, if \((x, y)\) tends to \((a, b)\) , then \((f(x), f(y))\) tends to \((f(a+), f(b-))\) by hypothesis; it suffices to apply formula (12) of II, p.~60 between x and y, and to pass to the limit to obtain (1).

Given a regulated function f on a non-compact interval \(I \subset R\) , with endpoints a and b (a \textless{} b), the condition for f to have an integral over I can be presented in the following manner. The compact intervals \(J \subset I\) form a directed set \(\mathfrak{K}(I)\) with respect to the relation \(\subset^{1}\) , for if \([\alpha, \beta]\) and \([\gamma, \delta]\) are two compact intervals contained in I, and if one puts \(\lambda = \min(\alpha, \gamma)\) , \(\mu = \max(\beta, \delta)\) , then the interval \([\lambda, \mu]\) is contained in I and contains the two intervals considered. For each compact interval \(J = [\alpha, \beta]\) contained in I, let us put

\[
\int_ {\mathrm{J}} \mathbf {f} (t) d t = \int_ {\alpha} ^ {\beta} \mathbf {f} (t) d t;
\]

for f to admit an integral over I it is necessary and sufficient that the map \(\mathbf{J} \mapsto \int_{\mathrm{J}} \mathbf{f}(t) dt\) have a limit with respect to the directed set \(\mathcal{K}(\mathrm{I})\) ; this limit is then the integral \(\int_{a}^{b} \mathbf{f}(t) dt\) , which we again denote by \(\int_{\mathrm{I}} \mathbf{f}(t) dt\) .

PROPOSITION 1 (Cauchy's criterion for integrals). Let \(\mathbf{f}\) be a regulated function on an interval \(\mathrm{I} \subset \mathbf{R}\) having endpoints \(a\) and \(b\) ( \(a < b\) ). For the integral \(\int_{a}^{b} \mathbf{f}(t) dt\) to exist it is necessary and sufficient that for every \(\varepsilon > 0\) there exist a compact interval \(\mathrm{J}_0 = [\alpha, \beta]\) contained in I, such that for any compact interval \(\mathrm{K} = [x, y]\) contained in I and having no interior points in common with \(\mathrm{J}_0\) , one has \(\left\| \int_{\mathrm{K}} \mathbf{f}(t) dt \right\| \leqslant \varepsilon\) .

Indeed, since E is complete the Cauchy criterion shows that for the integral \(\int_{\mathrm{I}} \mathbf{f}(t) dt\) to be convergent it is necessary and sufficient that for any \(\varepsilon > 0\) there exists a compact interval \(\mathrm{J}_0 = [\alpha, \beta]\) with for every compact interval J such that \(\mathrm{J}_0 \subset \mathrm{J} \subset I\) one has \(\left\| \int_{\mathrm{J}} \mathbf{f}(t) dt - \int_{\mathrm{J}_0} \mathbf{f}(t) dt \right\| \leqslant \varepsilon\) . The proposition will follow from the following lemma:

Lemma. Let \(J_{0} = [\alpha, \beta]\) be a compact interval contained in I. In order that \(\left\|\int_{J} \mathbf{f}(t) dt - \int_{J'} \mathbf{f}(t) dt\right\| \leqslant \varepsilon\) for every pair of compact intervals J, \(J'\) contained in I and containing \(J'\) it is necessary that \(\left\|\int_{K} \mathbf{f}(t) dt\right\| \leqslant \varepsilon\) , and it suffices that \(\left\|\int_{K} \mathbf{f}(t) dt\right\| \leqslant \varepsilon/2\) , for every compact interval K contained in I and having no interior point in common with \(J_{0}\) .

Indeed, if for \(J_{0} \subset J \subset I\) and \(J_{0} \subset J' \subset I\) , one has

\[
\left\| \int_ {J} \mathbf {f} (t) d t - \int_ {J ^ {\prime}} \mathbf {f} (t) d t \right\| \leqslant \varepsilon
\]

one sees in particular that, for \(x \leqslant y \leqslant \alpha\) , or for \(\beta \leqslant x \leqslant y\) (x and y in I), one has \(\left\|\int_{x}^{y} \mathbf{f}(t) dt\right\| \leqslant \varepsilon\) . Conversely, if \(\left\|\int_{K} \mathbf{f}(t) dt\right\| \leqslant \varepsilon/2\) for every compact interval \(K \subset I\) such that \(K \cap J_{0} = \emptyset\) , and if \(J = [x, y]\) , \(J' = [z, t]\) are two compact intervals contained in I and containing \(J_{0}\) , one has

\[
\left\| \int_ {J} \mathbf {f} (t) d t - \int_ {J ^ {\prime}} \mathbf {f} (t) d t \right\| = \left\| \int_ {x} ^ {z} \mathbf {f} (t) d t + \int_ {t} ^ {y} \mathbf {f} (t) d t \right\| \leqslant \varepsilon ,
\]

since

\[
x \leqslant \alpha \leqslant \beta \leqslant y \quad \text { and } \quad z \leqslant \alpha \leqslant \beta \leqslant t.
\]

Example. If the interval I is bounded, and if \(\mathbf{f}\) is bounded on I, then the integral \(\int_{\mathrm{I}}\mathbf{f}(t)dt\) always exists, for, by the mean value theorem, one has, for \(y\leqslant \alpha \leqslant \beta \leqslant z\)

\[
\left\| \int_ {y} ^ {\alpha} \mathbf {f} (t) d t \right\| \leqslant (\alpha - a) \sup _ {x \in I} \| \mathbf {f} (x) \|, \quad \left\| \int_ {\beta} ^ {z} \mathbf {f} (t) d t \right\| \leqslant (b - \beta) \sup _ {x \in I} \| \mathbf {f} (x) \|
\]

and it suffices to take \(\alpha - a\) and \(b - \beta\) small enough for the Cauchy criterion be satisfied. One may note that in this case a primitive of f on I does not necessarily have a right (resp. left) derivative at the left-hand endpoint (resp. right-hand endpoint) of I (when this number is finite) contrary to the situation when I is compact and f is regulated on I (cf.~II, p.~33, exerc. 1).

